\begin{figure}[tb]
  \centering
  %% reference
  \begin{subfigure}[t]{0.02\textwidth}
    % \vspace{.2cm}
    % \textbf{a}
    \raisebox{-1.25cm}{\rotatebox[origin=t]{90}{\textbf{Reference}}}
  \end{subfigure}
  \begin{subfigure}[t]{0.96\textwidth}
    \includegraphics[width=\linewidth,valign=t]{images/outlines_reference.png} \\
  \end{subfigure}\hfill
  %% Transformer
  \vspace{-.5cm}
  \begin{subfigure}[t]{0.02\textwidth}
    % \vspace{.2cm}
    % \textbf{b}
    \raisebox{-1.25cm}{\rotatebox[origin=t]{90}{\textbf{Transformer}}}
  \end{subfigure}
  \begin{subfigure}[t]{0.96\textwidth}
    \includegraphics[width=\linewidth,valign=t]{images/outlines_Transformer.png} \\
  \end{subfigure}\hfill
  %% single-step
  \vspace{-.5cm}
  \begin{subfigure}[t]{0.02\textwidth}
    % \vspace{.2cm}
    % \textbf{b}
    \raisebox{-1.25cm}{\rotatebox[origin=t]{90}{\textbf{Feedforward}}}
  \end{subfigure}
  \begin{subfigure}[t]{0.96\textwidth}
    \includegraphics[width=\linewidth,valign=t]{images/outlines_single_step.png} \\
  \end{subfigure}\hfill
  ~
    \caption{
    \textbf{(Prediction visualization)}
    A qualitative comparison of length-100 trajectories generated by the Trajectory Transformer and a feedforward Gaussian dynamics model from PETS, a state-of-the-art planning algorithm \cite{chua2018pets}.
    Both models were trained on trajectories collected by a single policy, for which a true trajectory is shown for reference.
    Compounding errors in the single-step model lead to physically implausible predictions, whereas the Transformer-generated trajectory is visually indistinguishable from those produced by the policy acting in the actual environment.
    The paths of the feet and head are traced through space for depiction of the movement between rendered frames.
    }
    \label{fig:humanoid}
    \vspace{-.3cm}
\end{figure}